\subsection{稠密性定理}

定理 3（雅各布森）. —— 设 M 是一个半单 A-模，c 是加群 M 的一个自同态。则 c 属于 M 的双交换子 \(A_{M}^{\prime\prime}\) 当且仅当它满足下列条件：

\begin{enumerate}
\def\labelenumi{(\Alph{enumi})}
\setcounter{enumi}{3}
\tightlist
\item
  对每个有限子集 \(\mathrm{F}\)（\(\mathrm{M}\) 的），存在一个元素 \(a\)（属于 \(\mathrm{A}\)），使得 \(c\) 与 \(a_{\mathrm{M}}\) 在 \(\mathrm{F}\) 上相合。
\end{enumerate}

首先设 \(c\) 满足条件 (D)。设 \(u\) 是 \(\mathrm{A}_{\mathrm{M}}^{\prime}\) 的一个元素。设 \(x\) 是 \(\mathrm{M}\) 的一个元素，对子集 \(\mathrm{F} = \{x, u(x)\}\) 应用条件 (D)。存在一个元素 \(a\)（属于 \(\mathrm{A}\)），使得 \(c(x) = ax\) 且 \(c(u(x)) = au(x)\)，于是 \(u(c(x)) = u(ax) = au(x) = c(u(x))\)。由于 \(x\) 是任意的，我们有 \(cu = uc\)；又因这对一切 \(u\) 成立，我们有 \(c \in \mathrm{A}_{\mathrm{M}}^{\prime \prime}\)。

为证其逆，我们使用下面的引理。

引理 3. —— 设 M 是一个半单 A-模。设 B 是 A-模 M 的双交换子。则 M 的每个 A-子模都是 M 的 B-子模。

设 N 是 M 的一个 A-子模。由 VIII, 第 56 页的推论 2，存在 A-模 M 的一个以 N 为像的投影算子 p；它属于 \(A_{M}^{\prime}\)。由于对每个 \(b \in B\) 我们有关系式 pb = bp，我们得到 N 是 M 的 B-子模。

现在来完成定理 3 的证明。设 c 属于 \(A_{M}^{\prime\prime}\)，且设 \(F = \{x_{1}, \ldots, x_{n}\}\) 是 M 的一个有限子集。把元素 \((x_{1}, \ldots, x_{n})\)（\(M^{n}\) 的）记作 x。A-模 \(M^{n}\) 是半单的，且由 VIII, 第 79 页的命题 2，它的双交换子与 \(A_{M}^{\prime\prime}\)-模 \(M^{n}\) 的位似重合。由引理 3，\(M^{n}\) 的 A-子模 Ax 是一个 \(A_{M}^{\prime\prime}\)-子模（\(M^{n}\) 的）。设 \(a \in A\) 使得 \((cx_{1}, \ldots, cx_{n})\) 等于 ax。于是 c 与 \(a_{M}\) 在 \(\{x_{1}, \ldots, x_{n}\}\) 上相合，这就蕴含条件 (D)。

注. —— 把 M 的加群的自同态环记作 E。给 M 赋以离散拓扑；环 E 由从 M 到 M 的映射组成，我们可以给它赋以由 \(M^{M}\) 上乘积拓扑诱导的拓扑（“M 中单收敛拓扑”，一般拓扑，I, §2, No.~3, 第 31 页）。E 上的拓扑是豪斯多夫的，且与 E 的加群结构相容。对 E 中每个 f，映射 \(g \mapsto f \circ g\) 与 \(g \mapsto g \circ f\)（从 E 到 E）都是连续的。因此，E 的每个子集的交换子在 E 中闭。于是定理 3 蕴含 \(A_{M}^{\prime\prime}\) 是 \(A_{M}\) 在 E 中的闭包。

